\chapter{A Concrete Naming Certificate for $\RegularCauchySubsequencePrincipleUp$}
\label{ch:concrete-instances-regular-cauchy-subsequence-principle-namecert}
\origin{human}

Following Bishop and Bridges on subsequences of regular Cauchy names, the $\RegularCauchySubsequencePrincipleUp$ packet records the finite route by which a source regular sequence may be inspected along a displayed monotone cofinal index map while preserving the rational readback and Real-facing seal. It sits beside the schedule rows of \autoref{ch:concrete-instances-streamname-namecert}, the regular rational carrier of \autoref{ch:concrete-instances-regseqrat-namecert}, the terminal real-name surface of \autoref{ch:concrete-instances-real-namecert}, the dyadic tolerance ledger of \autoref{ch:concrete-instances-dyadicratcore-namecert}, and the extractor route of \autoref{ch:concrete-instances-cauchy-subsequence-extractor-namecert}. The packet names subsequence stability as finite NameCert data; it is not quotient stream equality, not a selected limit, not countable choice, not host real equality, and not ambient completeness of $\RealUp$.

\begin{definition}[Regular Cauchy subsequence principle carrier]
\label{def:regular-cauchy-subsequence-principle-carrier}
A $\RegularCauchySubsequencePrincipleUp$ carrier is a finite $\BHist$ packet
$$
Q=(S,I,D,W,R,E,H,C,P,N)
$$
with the following displayed rows:
\begin{enumerate}
\item $S$ is the source regular sequence row, carried through the $\StreamNameUp$ and $\RegSeqRatUp$ schedules.
\item $I$ is the monotone cofinal index map, recording the finite request translation from subsequence windows back into source windows and the cofinal tail witness used by every downstream read.
\item $D$ is the $\DyadicRatCoreUp$ tail-budget ledger, carrying the displayed precision cells, dyadic slack comparisons, and tail thresholds.
\item $W$ is the finite $\StreamNameUp$ window row naming the source and subsequence observations inspected at each displayed precision request.
\item $R$ is the $\RegSeqRatUp$ readback row reached only after $S,I,D,W$ have been displayed.
\item $E$ is the terminal $\RealUp$ seal row, exposed only as the seal of the same cofinal regular readback route.
\item $H$ records componentwise $\hsame$ transport for $S,I,D,W,R,E,C,P,N$.
\item $C$ records displayed $\Cont$ replay from a subsequence precision request through $I$, the source row $S$, the dyadic tail budget $D$, the finite windows $W$, the regular readback $R$, and the Real seal $E$.
\item $P$ is the $\Pkg$ provenance over $\RegularCauchySubsequencePrincipleUp$, $\StreamNameUp$, $\RegSeqRatUp$, $\RealUp$, $\DyadicRatCoreUp$, $\BHist$, $\hsame$, $\Cont$, and $\NameCert$ rows.
\item $N$ is the local $\NameCert_{\RegularCauchySubsequencePrincipleUp}$ row.
\end{enumerate}
The classifier compares accepted packets only through $S,I,D,W,R,E,H,C,P,N$. It has no coordinate for quotient stream equality, selected subsequence limits, countable choice, host real equality, or ambient $\RealUp$ completeness.
\end{definition}

\begin{theorem}[Regular Cauchy subsequence principle NameCert obligations]
\label{thm:regular-cauchy-subsequence-principle-namecert-obligations}
For every accepted $\RegularCauchySubsequencePrincipleUp$ carrier $Q=(S,I,D,W,R,E,H,C,P,N)$, the local $\NameCert_{\RegularCauchySubsequencePrincipleUp}$ surface consists exactly of carrier admission, cofinal index stability, regular tail preservation, Real-seal handoff, componentwise $\hsame$ transport, displayed $\Cont$ replay, $\Pkg$ provenance, and local naming. Every subsequence read must expose the source regular sequence, monotone cofinal index map, dyadic tail budget, finite StreamName windows, and RegSeqRat readback before the Real seal is consumed.
\end{theorem}
\begin{proof}
Project the carrier of \autoref{def:regular-cauchy-subsequence-principle-carrier}. The analytic rows are $S,I,D,W,R,E$, and the structural rows $H,C,P,N$ only transport, replay, source, and name those rows.

The cofinal index row $I$ is the sole route from a subsequence precision request back to the source regular sequence. The dyadic budget $D$ and finite window row $W$ are read on that route, so regular tail preservation is tied to the displayed source sequence rather than to a detached subsequence object.

The readback row $R$ and Real seal $E$ are downstream of the same finite route. Since the packet contains no quotient coordinate, selected limit row, choice schedule, host equality row, or ambient completeness theorem, the listed obligations exhaust the local NameCert surface.
\end{proof}

\begin{theorem}[Regular Cauchy subsequence Real-seal handoff]
\label{thm:regular-cauchy-subsequence-principle-real-seal-handoff}
Every $\RealUp$-facing read of an accepted $\RegularCauchySubsequencePrincipleUp$ carrier factors through the displayed cofinal tail and regular-readback route
$$
I,S,D,W \longmapsto R \longmapsto E .
$$
Thus a downstream Real consumer may cite the seal row $E$ only after the monotone cofinal index map, source regular sequence, dyadic tail budget, finite StreamName windows, and RegSeqRat readback have all been displayed. The handoff exports no quotient stream, selected limit, countable-choice schedule, host real equality, or ambient Real completeness.
\end{theorem}
\begin{proof}
Start with a consumer asking for $E$. The carrier definition places $E$ after the readback row $R$, and $R$ is admitted only after the cofinal index, source sequence, dyadic budget, and finite windows have been displayed.

The continuation row $C$ records the same order of exposure. A subsequence request is replayed through $I$ before it reads $S$, then through $D$ and $W$ before it reaches the regular rational row.

Transport and provenance through $H$ and $P$, with local naming through $N$, preserve this finite packet. Therefore every Real-facing read is mediated by the displayed cofinal tail route and receives none of the excluded quotient, choice, equality, or completeness data.
\end{proof}

\closureat{RegularCauchySubsequencePrincipleUp}{seedStr}

\begin{closurestatus}{\RegularCauchySubsequencePrincipleUp}
  \constructivestory{A $\RegularCauchySubsequencePrincipleUp$ packet is generated from a source regular sequence, a monotone cofinal index map, a DyadicRatCore tail budget, finite StreamName windows, RegSeqRat readback, a RealUp seal, componentwise hsame transport, displayed Cont replay, Pkg provenance, and a local NameCert row.}
  \theoryclosure{\seedClosure}
  \scopeclosed{\autoref{def:regular-cauchy-subsequence-principle-carrier}, \autoref{thm:regular-cauchy-subsequence-principle-namecert-obligations}, and \autoref{thm:regular-cauchy-subsequence-principle-real-seal-handoff} seed the finite cofinal subsequence route from displayed regular Cauchy data to the downstream $\RealUp$ seal.}
  \formalstatus{\unformalizedV}
  \bridgestatus{none}
  \notclaimed{This chapter does not close quotient stream equality, selected subsequence limits, countable choice, host real equality, ambient $\RealUp$ completeness, Lean verification, or bridge checking.}
  \upgradepath{Obligation closure requires checked carrier admission, cofinal index stability, dyadic tail-budget exactness, finite-window replay, regular readback, Real seal ordering, transport, replay, provenance, and local naming for BEDC.Derived.RegularCauchySubsequencePrincipleUp.}
\end{closurestatus}
